\chapter*{Introduction}
Hand exoskeletons have emerged as promising devices for motor rehabilitation, assisting patients in recovering finger mobility after neurological injury or surgery. By guiding or assisting the natural motion of the hand, these devices enable repetitive, task-specific exercises that are known to promote neuroplasticity and functional recovery. However, the development of hand exoskeletons poses significant engineering challenges. The human hand is a mechanically complex structure, and any device that interacts with it must satisfy stringent requirements in terms of kinematic compatibility, force compliance, and — critically — safety. A malfunction that applies an unintended force to an impaired hand can cause pain, injury, or loss of trust in the rehabilitation process.
 
Traditional approaches to the development of such devices rely heavily on iterative hardware prototyping: a mechanism is designed, built, instrumented, tested, redesigned, and tested again. This cycle is effective but slow and expensive, and it exposes both the device and the user to the risks inherent in testing immature hardware. Moreover, early-stage design decisions — such as the choice of link dimensions, actuator sizing, or control gains — are difficult to evaluate without a working prototype, leading to suboptimal designs that must be corrected in later iterations.
 
The Digital Twin paradigm offers an alternative. By constructing a virtual representation of the physical system that reproduces its mechanical, dynamic, and control behavior with sufficient fidelity, the designer can evaluate design choices, tune control parameters, and verify safety strategies entirely in simulation, before any hardware is built or any patient is involved. The Digital Twin does not replace hardware testing, but it can dramatically reduce the number of physical iterations required and can identify critical issues — such as instabilities, excessive forces, or undetected fault conditions — at a stage where they are inexpensive to correct.
 
This thesis develops a comprehensive Digital Twin for a hand exoskeleton, covering the full pipeline from mechanical design to functional safety verification. The work builds upon the exoskeleton design presented in Bartalucci's thesis [1], from which the geometric dimensions of the finger mechanism are derived. The Digital Twin is implemented entirely within the ROS 2 framework, using the Pinocchio rigid-body dynamics library [2] for kinematic and dynamic computations, and is organized as a modular workspace of nine ROS 2 packages that can be independently developed, tested, and deployed.
 
The exoskeleton mechanism is characterized by multiple closed kinematic chains — a common feature in finger exoskeletons, where parallel linkages are used to enforce coordinated motion among the phalanges. Because the standard URDF robot description format does not support closed-loop topologies, the mechanism was modeled as an equivalent open-chain structure in which the closed loops are opened at selected joints and the loop-closure conditions are reimposed numerically at runtime. This hybrid approach — a URDF-compatible tree structure augmented by software-enforced geometric constraints — enables the use of standard robotic tools while preserving the physical consistency of the original design.
 
On this kinematic foundation, a reduced-order dynamic model was developed by exploiting the single-degree-of-freedom nature of the mechanism. The full multibody dynamics, including inertial, Coriolis, gravitational, friction, and passive biomechanical effects, is projected onto the actuated crank coordinate through a constraint-consistent reduction that yields a scalar nonlinear equation. This equation is evaluated at each simulation step using the quantities computed by the Pinocchio library, avoiding any offline symbolic reduction and ensuring that the model remains consistent with the URDF-based representation at all times.
 
A cascade admittance-trajectory control architecture was implemented to regulate the interaction between the user and the exoskeleton. The outer loop encodes the desired interaction behavior through virtual dynamics — a configurable mass-damper-spring system whose parameters shape the perceived compliance of the device — while the inner loop ensures accurate trajectory tracking through a PD controller with model-based feed-forward compensation. The cascade structure decouples the interaction design from the low-level actuator regulation, enabling independent tuning of the two loops.
 
A central contribution of this thesis is the development of a functional safety architecture that addresses the specific risks of a device intended for direct physical human-robot interaction. The architecture follows a defense-in-depth philosophy, distributing safety responsibilities across three cooperating layers: an enforcement gateway that physically limits the signals reaching the actuator, a fault detection plugin that monitors the system against configurable thresholds, and a governance state machine that manages the safety state through a well-defined transition graph with fault latching and bridge supervision. The entire safety subsystem is built on a plugin-based architecture that allows new safety modes to be added without modifying the core framework. A fault injection framework, capable of introducing six types of controlled faults on any of the four critical signal channels, was developed to enable systematic validation of the safety mechanisms.
 
Complementing the threshold-based safety layer, a model-based fault detection module was developed to explore the feasibility of observer-based sensor health monitoring. Four complementary observer layers — exploiting different physical principles and operating at different detection timescales — were implemented and integrated into the Digital Twin. While this module was not the primary focus of the thesis and has not yet been connected to the safety decision chain, it constitutes a functional proof of concept whose integration path is clearly defined within the existing architecture.
 
Finally, the hardware prototype of the hand exoskeleton was documented in detail, covering the electromechanical components, the wiring architecture, the motor drive system, and the board-level software modifications introduced to support the Digital Twin integration. This documentation establishes the complete interface specification for a future hardware-in-the-loop deployment in which the ROS 2 stack running on a workstation communicates with the physical prototype through a WebSocket-based distributed control architecture.
 
The thesis is organized as follows. Chapter 1 describes the mechanical design of the exoskeleton components in Fusion 360 and the generation of the URDF model, including the handling of closed kinematic chains. Chapter 2 presents the kinematic modeling and the constrained forward kinematics formulation. Chapter 3 develops the reduced-order dynamic model and validates it in open-loop conditions. Chapter 4 describes the cascade admittance-trajectory control architecture and its closed-loop validation. Chapter 5 presents the functional safety architecture, including the state machine, the enforcement bridge, the fault detection plugin, and the fault injection framework. Chapter 6 describes the model-based fault detection module and its four observer layers. Chapter 7 documents the hardware prototype of the hand exoskeleton. Chapter 8 draws the conclusions and outlines directions for future work.